% Generated draft: requires mathematical and source review.
\documentclass[11pt]{article}
\usepackage[margin=25mm]{geometry}
\usepackage{fontspec}
\setmainfont{Latin Modern Roman}
\usepackage{amsmath,amssymb,mathtools,unicode-math}
\setmathfont{Latin Modern Math}
\usepackage{enumitem,needspace,xurl,hyperref}
\setlength{\parindent}{0pt}
\setlength{\parskip}{5pt}
\setlength{\emergencystretch}{3em}
\newcommand{\N}{\mathbb{N}}
\newcommand{\Z}{\mathbb{Z}}
\newcommand{\Q}{\mathbb{Q}}
\newcommand{\R}{\mathbb{R}}
\newcommand{\C}{\mathbb{C}}
\newcommand{\F}{\mathbb{F}}
\newcommand{\A}{\mathbb{A}}
\newcommand{\PP}{\mathbb{P}}
\newcommand{\E}{\mathbb{E}}
\newcommand{\eps}{\varepsilon}
\newcommand{\dd}{\,\mathrm{d}}
\begin{document}
\section*{UnsolvedMath ID: 10400136}
\textbf{Source problem code:} catalogue entry
\subsection{Definitions and mathematical statement}
For an admissible positive integer \(N\), write
\[\mu_N=\{\zeta\in\C:\zeta^N=1\}.\]
The state sum \(H_N(T)\) in the source has a phase ambiguity given by multiplication by elements of \(\mu_N\). For a nonzero complex representative \(z\), this means that the available phase-independent data determine at most its orbit \(z\mu_N\), unless further normalization is supplied. The problem attributed to S. Baseilhac and R. Benedetti first asks to understand this ambiguity: determine the actual root-of-unity factors arising from changes of the source's auxiliary data \(T\) representing \((W,L,\rho)\), including how these factors depend on the choices and transformations. The excerpt does not define \(T\), the state sum or the transformations, so no specific phase formula is supplied here.

The second, explicitly tentative request is to use this analysis to obtain an invariant of \((W,L,\rho)\) equipped with extra structure, refining \(K_N(W,L,\rho)\). In precise structural terms, one would specify the extra datum \(s\), its equivalence/transport rules and an invariant \(\widetilde K_N(W,L,\rho,s)\) whose prescribed forgetting operation recovers \(K_N(W,L,\rho)\), while retaining phase information discarded there. Neither the nature of \(s\), the refined target nor the actual relation between \(H_N\) and \(K_N\) is stated. In particular this draft does not identify \(K_N\) with a root-of-unity orbit or with a chosen power of \(H_N\); those would require the missing definitions. The words ``possibly derive'' are retained as a proposed application, not an assertion that a canonical phase lift exists.
\subsection{Short English statement}
Determine the root-of-unity ambiguity in the state sum \(H_N(T)\). Baseilhac and Benedetti also suggest using suitable extra structure to retain this phase and refine the invariant \(K_N\).
\subsection{Sources}
\begin{itemize}
\item Local UnsolvedMath registry, permanent ID \(10400136\), exact source code \texttt{catalogue entry}.
\item S. Baseilhac and R. Benedetti, attribution in Ohtsuki, \emph{Problems on invariants of knots and 3-manifolds} (2002), Problem 7.21, PDF page 114: \url{https://msp.org/gtm/2002/04/gtm-2002-04-024s.pdf}. Registry access 2026-07-29, pdf-font-structure extraction.
\end{itemize}
\end{document}
